\section{Theoretical Framework and Approach}
\subsection{Approach}

This is a conceptual article. It builds a model from existing theory and states its claims as propositions, which are statements about relationships among constructs that are meant to be tested rather than claims that have already been confirmed. Among the designs available for conceptual work, ours is closest to what \citet{jaakkola2020designing} calls a model paper, one that specifies constructs and the relationships among them. We follow the view of \citet{whetten1989constitutes} that a theoretical contribution must say which factors matter, how they are related, and why the proposed relationships should hold, and we adopt the propositional style that \citet{cornelissen2017developing} describes as suited to explaining variance in an outcome across conditions. Throughout, we separate three kinds of statement. The first reports what cited studies found. The second reports what cited theories argue. The third, flagged each time, is our own reasoning, which has not yet been tested.

\subsection{Psychological Contract Theory and Ideological Currency}

A psychological contract is an individual's belief about the reciprocal obligations between that individual and the employing organization \citep{rousseau1989psychological}. The belief need not match what the organization intended, and it need not be written down. What makes it a contract, in the psychological sense, is that the employee believes a promise was made and that something was given in return. \citet{rousseau2001schema} argued that these beliefs form partly before entry, from schemas the person already holds about employment and from promises communicated through words and actions during recruitment and early socialization. That point about formation matters for our model, because it treats where and how an obligation was communicated as part of the contract itself.

\citet{thompson2003violations} extended the theory by adding ideological currency to the transactional and relational currencies already recognized. An ideological contract involves the employee's belief that the organization has made a credible commitment to pursue a valued cause or principle, one that benefits parties beyond the employee. When such a commitment is broken, the employee has lost more than a personal benefit, because a principle the employee cared about has been abandoned. This is the closest established construct to what we mean by an employer's ethical commitments, and it is why we ground the model here rather than in a general account of fairness. Psychological contract theory also supplies the distinction between breach, the cognitive judgment that an obligation has gone unmet, and violation, the emotional experience of anger and betrayal that may or may not follow \citep{morrison1997employees}. That distinction lets us ask what turns a noticed breach into a felt violation.

\subsection{The Exit-Voice-Loyalty Framework and the Question of Venue}

Hirschman's framework concerns how members respond when an organization they belong to deteriorates. They can leave, which is exit, or they can try to change it from within, which is voice; loyalty holds exit back and gives voice time to work \citep{hirschman1970exit}. Organizational researchers later added neglect and arranged the responses along two dimensions, active versus passive and constructive versus destructive \citep{farrell1983exit}, and showed that satisfaction, investment, and the quality of alternatives shape which response people choose \citep{rusbult1988impact}. \citet{turnley1999impact} connected this tradition directly to psychological contracts, reporting that managers who experienced contract violations reported more exit, voice, and neglect and less loyalty.

The framework is a theory of response choice after something has gone wrong, which is exactly the problem we address. Its limitation for our purpose lies in how voice has been carried into organizational research. Hirschman's own notion of voice was broad: any attempt to change a state of affairs rather than escape it, whether by petitioning management, appealing to a higher authority, or protesting in ways meant to rally outside opinion. Organizational behavior research later narrowed the construct, typically defining voice as upward communication of concerns or suggestions to someone with authority inside the organization \citep{morrison2014employee}, while disclosures to outside parties were studied separately as external whistleblowing. The result is that neither literature asks what makes an employee choose one venue over the other after the same breach. A complaint raised with a supervisor and a complaint posted to a public audience differ in who hears them and in the risk they carry for the employee. We therefore keep the framework's core logic, in which the availability and expected effectiveness of voice govern whether people stay and speak or leave, and we treat venue as a choice within voice that needs its own explanation.

\subsection{Why These Two Theories}

Several alternatives were considered and set aside. Signaling theory has been used to explain why job seekers are drawn to socially responsible employers \citep{jones2014why}, and it has been combined with ideological contract reasoning to study online whistleblowing about corporate social advocacy \citep{yang2026managing}. It is well suited to explaining how public claims create expectations, but it says little about what employees do after those expectations are disappointed, which is our focal question. Social exchange reasoning in general underlies psychological contract theory, but on its own it is too broad to generate predictions about channel choice. Generational theory, traced to \citet{mannheim1952problem}, is not used here as a causal theory, for reasons developed in the literature review: its workplace applications have weak empirical support, and we prefer to treat the generational question as something to be tested against a competing explanation. Psychological contract theory tells us what an ethical promise is and why its breach can be felt as betrayal. The exit-voice-loyalty framework tells us how people choose among responses. The two fit together because the second begins where the first ends.

\subsection{Scope Conditions}

The model applies to employees of organizations that make ethical commitments that employees can perceive, whether through public statements, recruitment messages, or interpersonal assurances. It concerns breaches of those commitments as the employee perceives them; it does not assume that a breach occurred in any objective sense. It is not a model of whistleblowing on crime in general, where legal duties and protections change the calculus, although it overlaps with that literature. It also assumes that public channels are available to the employee at a cost that is not prohibitive; where public criticism of an employer carries severe legal or social penalties, we expect the model's predictions about public voice to weaken, a point we return to in the discussion.
